\sectionkicker{Source-backed technical notes}
\subsection*{AgentはChatを出る — Harness・Context・Computer/Terminal Use: Technical Notes}
本欄は記事本文で圧縮した一次資料上の情報を、比較・再検証しやすい形で整理したものである。時系列、確認済みの主張、留意点、一次資料URLを掲載する。
\medskip
\subsection*{Theme at a glance}
\addcontentsline{toc}{subsection}{Theme at a glance}
\begin{center}
\footnotesize
\begin{tabularx}{\linewidth}{@{}p{0.20\linewidth}p{0.10\linewidth}p{0.14\linewidth}X@{}}
\toprule
資料 & 位置づけ & 種別 & 時系列 \\
\midrule
OpenAI January agent-system engineering: Codex loop and in-house data agent & 主要資料 & Agent & 2026-01-23 (技術解説); 2026-01-29 (技術解説) \\
Claude Cowork research preview and new Claude constitution & 主要資料 & Agent & 2026-01-12 (研究Preview); 2026-01-22 (モデル行動方針公開) \\
Gemini API January 2026 tool and lifecycle updates & 主要資料 & API & 2026-01-29 (APIツール公開); 2026-01-12 (APIライフサイクル機能); 2026-01-08 (API入力拡張) \\
Mistral Vibe 2.0 & 主要資料 & Agent & 2026-01-27 (製品公開) \\
\bottomrule
\end{tabularx}
\end{center}
\normalsize

\begin{technicalnote}{OpenAI January agent-system engineering: Codex loop and in-house data agent}{主要資料}
\begingroup
% reader-facing Technical Notes break policy
\widowpenalty=10000
\clubpenalty=10000
\displaywidowpenalty=10000
\begin{tabularx}{\linewidth}{@{}>{\bfseries}p{0.22\linewidth}X@{}}
組織 & OpenAI \\
種別 & Agent \\
時系列 & 2026-01-23 (技術解説); 2026-01-29 (技術解説) \\
\end{tabularx}
\smallskip
{\bfseries 一次資料から整理したtechnical points}
\begin{itemize}[leftmargin=1.5em,itemsep=0.35em]
\item \textbf{一次情報で確認できる事実}: OpenAIのCodex記事はagent loopを、user、model inference、tools、conversation history、context-window managementを編成するharness logicとして説明する。
\item \textbf{一次情報で確認できる事実}: OpenAIのin-house data-agent記事は、metadata、human annotation、code enrichment、institutional knowledge、memory、runtime contextを重ね、pass-through access controlを用いるinternal-only GPT-5.2 agentを説明する。
\item \textbf{Vendor claim}: data agentはpublic productではなくinternal toolであり、分析高速化やreliability向上の例はOpenAIの運用経験で、独立benchmarkではない。
\end{itemize}
\begin{minipage}{\linewidth}
% reader-facing Technical Notes generic boundary/source tail group
{\bfseries 読む際の境界}
\begin{itemize}[leftmargin=1.5em,itemsep=0.35em]
\item \textbf{分析上の留意点}: 二つのsourceは異なるsystemを説明しており、一つのarchitecture clusterで扱う場合も帰属を分ける必要がある。
\item \textbf{分析上の留意点}: in-house data agentの知見は組織内case studyであり、普遍的なarchitecture benchmarkへ一般化しない。
\end{itemize}
\begin{samepage}
{\bfseries 一次資料}
\begingroup\sloppy
\Urlmuskip=0mu plus 2mu\relax
\begin{itemize}[leftmargin=1.5em,itemsep=0.25em]
\item {\scriptsize\url{https://openai.com/index/inside-our-in-house-data-agent}}
\item {\scriptsize\url{https://openai.com/index/unrolling-the-codex-agent-loop}}
\end{itemize}
\endgroup
\end{samepage}
\end{minipage}
\endgroup
\end{technicalnote}
\medskip

\begin{technicalnote}{Claude Cowork research preview and new Claude constitution}{主要資料}
\begingroup
% reader-facing Technical Notes break policy
\widowpenalty=10000
\clubpenalty=10000
\displaywidowpenalty=10000
\begin{tabularx}{\linewidth}{@{}>{\bfseries}p{0.22\linewidth}X@{}}
組織 & Anthropic \\
種別 & Agent \\
時系列 & 2026-01-12 (研究Preview); 2026-01-22 (モデル行動方針公開) \\
\end{tabularx}
\smallskip
{\bfseries 一次資料から整理したtechnical points}
\begin{itemize}[leftmargin=1.5em,itemsep=0.35em]
\item \textbf{一次情報で確認できる事実}: Anthropicの1月13日Labs announcementは、前日にCoworkがresearch previewとして公開され、Claudeのagentic capabilityをdesktop workflowへ持ち込んだと記載する。
\item \textbf{一次情報で確認できる事実}: Anthropicは1月22日にClaudeの新constitutionを公開し、model training processの重要な一部としてClaudeのbehaviorを直接方向付ける文書だと説明している。
\end{itemize}
\begin{minipage}{\linewidth}
% reader-facing Technical Notes generic boundary/source tail group
{\bfseries 読む際の境界}
\begin{itemize}[leftmargin=1.5em,itemsep=0.35em]
\item \textbf{分析上の留意点}: Coworkは発表時点でresearch previewであり、その時点のproduct boundaryを後日のavailabilityへ一般化しない。
\item \textbf{分析上の留意点}: constitutionは意図されたbehaviorとtraining guidanceを示すが、Anthropic自身が述べる通り、実際のmodel outputが常にその理想に従うとは限らない。
\end{itemize}
\begin{samepage}
{\bfseries 一次資料}
\begingroup\sloppy
\Urlmuskip=0mu plus 2mu\relax
\begin{itemize}[leftmargin=1.5em,itemsep=0.25em]
\item {\scriptsize\url{https://www.anthropic.com/news}}
\end{itemize}
\endgroup
\end{samepage}
\end{minipage}
\endgroup
\end{technicalnote}
\medskip

\begin{technicalnote}{Gemini API January 2026 tool and lifecycle updates}{主要資料}
\begingroup
% reader-facing Technical Notes break policy
\widowpenalty=10000
\clubpenalty=10000
\displaywidowpenalty=10000
\begin{tabularx}{\linewidth}{@{}>{\bfseries}p{0.22\linewidth}X@{}}
組織 & Google \\
種別 & API \\
時系列 & 2026-01-29 (APIツール公開); 2026-01-12 (APIライフサイクル機能); 2026-01-08 (API入力拡張) \\
\end{tabularx}
\smallskip
{\bfseries 一次資料から整理したtechnical points}
\begin{itemize}[leftmargin=1.5em,itemsep=0.35em]
\item \textbf{一次情報で確認できる事実}: Gemini API changelogは1月29日、gemini-3-pro-previewとgemini-3-flash-previewへのComputer Use tool supportを記録している。
\item \textbf{一次情報で確認できる事実}: changelogは1月12日のmodel lifecycle機能と、1月8日のCloud Storageおよびpublic/private pre-signed database URL入力対応（100 MB file limit）を記録する。
\item \textbf{一次情報で確認できる事実}: 同changelogには1月中のVeo 4K／portrait-output変更と、複数のalias・deprecation eventも記録されている。
\end{itemize}
\begin{minipage}{\linewidth}
% reader-facing Technical Notes generic boundary/source tail group
{\bfseries 読む際の境界}
\begin{itemize}[leftmargin=1.5em,itemsep=0.35em]
\item \textbf{分析上の留意点}: Computer Use supportはpreview model identifierに紐付くため、発表時点のpreview／deprecation statusを保持して扱う。
\item \textbf{分析上の留意点}: changelogが確認するのはAPI availabilityであり、computer-use behaviorの品質や安全性を独立評価する資料ではない。
\end{itemize}
\begin{samepage}
{\bfseries 一次資料}
\begingroup\sloppy
\Urlmuskip=0mu plus 2mu\relax
\begin{itemize}[leftmargin=1.5em,itemsep=0.25em]
\item {\scriptsize\url{https://ai.google.dev/gemini-api/docs/changelog}}
\end{itemize}
\endgroup
\end{samepage}
\end{minipage}
\endgroup
\end{technicalnote}
\medskip

\begin{technicalnote}{Mistral Vibe 2.0}{主要資料}
\begingroup
% reader-facing Technical Notes break policy
\widowpenalty=10000
\clubpenalty=10000
\displaywidowpenalty=10000
\begin{tabularx}{\linewidth}{@{}>{\bfseries}p{0.22\linewidth}X@{}}
組織 & Mistral AI \\
種別 & Agent \\
時系列 & 2026-01-27 (製品公開) \\
\end{tabularx}
\smallskip
{\bfseries 一次資料から整理したtechnical points}
\begin{itemize}[leftmargin=1.5em,itemsep=0.35em]
\item \textbf{一次情報で確認できる事実}: Mistralは1月27日、Devstral 2 model familyを用いるterminal-native coding agentとしてVibe 2.0を公開した。
\item \textbf{一次情報で確認できる事実}: releaseはcustom subagents、multi-choice clarification、slash-command skills、unified agent modes、automatic updatesを追加し、Le Chat Pro/TeamとPAYGまたはBYOK経路で提供された。
\end{itemize}
\begin{minipage}{\linewidth}
% reader-facing Technical Notes generic boundary/source tail group
{\bfseries 読む際の境界}
\begin{itemize}[leftmargin=1.5em,itemsep=0.35em]
\item \textbf{分析上の留意点}: announcementはvendor authoredであり、生産性やmodel qualityに関する主張は独立再現していない。
\item \textbf{分析上の留意点}: Devstral 2自体は1月以前から存在するため、月次eventは新しい基盤model familyではなくVibe 2.0のagent／harness更新である。
\end{itemize}
\begin{samepage}
{\bfseries 一次資料}
\begingroup\sloppy
\Urlmuskip=0mu plus 2mu\relax
\begin{itemize}[leftmargin=1.5em,itemsep=0.25em]
\item {\scriptsize\url{https://mistral.ai/news/}}
\end{itemize}
\endgroup
\end{samepage}
\end{minipage}
\endgroup
\end{technicalnote}
\medskip
